% sections/s1_intro.tex -- Introduction.
% Body only; uses macros.tex; all citation keys are in refs.bib.
% The introduction states results in words and points to the section that establishes each of them; the numbers
% are given there.  The few numbers kept here carry a % src comment.
% Paths in "% src:" comments are relative to the root of the repository (see README.md).
\section{Introduction}\label{s1_intro:sec}

The time a confined random walker needs to reach a target, the first-passage time, is the basic quantity of
search and reaction problems \citep{Redner2001,BenichouVoituriez2014}. Its distribution is usually summarised by
its mean, but in bounded domains the mean need not be representative: the most probable value can lie far below
it \citep{Mattos2012,GodecMetzler2016,Grebenkov2018}. We study this gap in the simplest lattice setting. A lazy
walker attempts, at each step and with probability $q$, a move to one of the $2\dd$ nearest neighbours on the
hypercubic lattice $\Om=\{0,\dots,N-1\}^{\dd}$; a move that would leave the lattice is cancelled, and one site is
absorbing. We compare the mean first-passage time (MFPT) $\MF$ with the mode $\tmode$, the most probable
first-passage time, mainly from one corner to the opposite corner, and ask how both depend on the size $N$, on
the dimension $\dd=1,2,3$ and, on the square lattice, on a fraction $p$ of blocked sites (inert defects).

\paragraph{What is known.}
The generating-function formalism of \citet{MontrollWeiss1965} expresses the first-passage law through the
propagator of the walk without target; with it \citet{Montroll1969} found that the mean time to a single trap on
a periodic lattice of $\nn$ sites grows as $\nn\ln\nn$ in two dimensions and as $\nn$ in three, and the survival
probability of a walk on a finite lattice with a single trap was studied by \citet{WeissHavlinBunde1985}.
\citet{Giuggioli2020} gave the exact propagator of the lazy walk in reflecting hypercubes of any dimension and,
from it, the MFPT as a nested finite sum; for the rectangle, with the same boundary rule, the cosine double sum was
given for this walk by \citet{CondaminBenichou2006} (in resistance form already by \citet{Wu2004}), and pseudo-Green-function methods give MFPTs in general confined domains
\citep{Condamin2005,Condamin2007PRE,Condamin2007Nature}. By the commute-time identity
\citep{Chandra1989,Tetali1991} mean hitting times of a symmetric walk are determined by effective resistances, and
the resistance between two nodes of a rectangular grid of resistors with free boundaries is known as a double sum
\citep{Wu2004} and as a single sum for two arbitrary nodes and any grid size \citep{TanTan2020}; the
corner-to-corner resistance has a large-size expansion \citep{EssamWu2009} whose coefficients are known for every
aspect ratio \citep{IzmailianHuang2010}.

The distribution in confinement is also known in outline. Rescaled by the mean it is close to exponential when
exploration is non-compact \citep{Benichou2010}, and for large volumes the rescaled laws fall into universality classes
\citep{Meyer2011}; direct trajectories set a short time scale and
trajectories that explore the whole domain a much longer one \citep{GodecMetzler2016}; and for reversible Markov
chains the exponential approximation of hitting times from a stationary start has explicit error bounds
\citep{AldousBrown1992,AldousFill}, and relaxation and first-passage time scales interlace
\citep{HartichGodec2018,HartichGodec2019}. For diffusion to a small target the slow decay rate follows from the
perturbation of eigenvalues by a small removed region \citep{WardKeller1993}, the mean time to a small target on a
reflecting boundary has an extensive asymptotic theory \citep{HolcmanSchuss2014}, and uniform
approximations of the first-passage density combine that long-time behaviour with a short-time correction
\citep{IsaacsonNewby2013}; the same authors observe that the most probable time, unlike the mean, depends strongly
on the starting position, and compute it numerically. \citet{GodecMetzler2016} and \citet{Grebenkov2018} give the
most probable time of direct trajectories, and \citet{Grebenkov2018} relate the median to the mean.
\citet{SarvaharmanGiuggioli2023} gave an analytic framework for lattice walks with inert heterogeneities and showed, for
a single barrier on a chain, that its position decides whether the mean or only the mode and the tail respond; the
reduction of conductivity and diffusivity on site-diluted
lattices is a classical subject \citep{Kirkpatrick1973,WatsonLeath1974,Nieuwenhuizen1986,Ernst1987,HavlinBenAvraham1987}.

\paragraph{Contributions.}
Each item states its status --- proved (by hand or with computer assistance), computed exactly at the sizes stated
(with no claim beyond them), or derived by formal asymptotics --- and what was already known.

\begingroup
\renewcommand{\labelenumi}{(\roman{enumi})}
\begin{enumerate}
\item \emph{Means (proved; known in substance).} A cosine sum for the MFPT between any two sites in any dimension
(Theorem~\ref{s4_mfpt:thm-cosine}); in two dimensions single sums for opposite corners, in three equivalent forms,
and for an arbitrary pair (Theorems~\ref{s4_mfpt:thm-single} and~\ref{s4_mfpt:thm-pair}), valid for every $N\ge2$;
the large-$N$ expansion with closed-form coefficients (Theorem~\ref{s4_mfpt:thm-expansion}) and, with computer
assistance, an explicit remainder for every $N$ (Theorem~\ref{s4_mfpt:thm-remainder}); in three dimensions a double
sum and, with computer assistance, $q\MF=4.32046\ldots\,N^{3}-3.887\ldots\,N^{2}+\Order(1)$, the leading constant
being a sum of lattice Green's function values (Theorem~\ref{s4_mfpt:thm-3d}).
% src: data/msc_rigorous_spectral/06_certified_constants.json (C_3, c_tau_3D); data/msc_rigorous_spectral/16_certified_d3_remainder.json
Through the commute-time identity the two-dimensional formulas are resistor-grid results in print
(Section~\ref{s4_mfpt:ssec-lit}); we add a self-contained proof in the random-walk setting and expressions of three
expansion coefficients through the lemniscate constant, and record an overflow-free form of the single sum, an
elementary rearrangement that is the one to use in floating point.

\item \emph{Distributions and modes (exact computation).} First-passage distributions are computed by time
stepping, by a closed form in one dimension and by Laplace inversion, and checked by a separately written pole expansion;
a Cauchy--Schwarz bound, together with exact certificates for corner-to-corner transport at $q\in\{\frac12,\frac45,1\}$
and an a-priori bound on round-off otherwise, certifies that every discrete-time mode obtained by time stepping on the
lattice without defects is the global maximiser (Section~\ref{s3_methods:sec}); on lattices with defects the same
certificate is evaluated in double precision (Section~\ref{s7_defects:sec}). In one dimension $\tmode/\MF$ tends to a constant; in two dimensions $q\tmode/N^{2}$
increases at every computed size while $\tmode/\MF$ falls; in three dimensions $\tmode$ follows $N^{2}\ln N$ whereas
$\MF$ grows as $N^{3}$ (Section~\ref{s5_modes:sec}). That the mode lies below the mean in confinement is known; the
exact lattice values and their dependence on $N$ and $\dd$ are computed here.

\item \emph{A closed-form approximation of the mode (formal derivation).} Two-pole asymptotics of the first-passage
transform give a parameter-free formula for the mode in terms of $X=\muone q\MF$, where $\muone$ is the slowest
relaxation rate of the reflecting lattice (Formal result~\ref{s6_mechanism:res-L1}); its accuracy is measured
against the exact modes, and it fails in one dimension and for some placements of start and target. Its
ingredients are the lattice counterparts of known continuum results \citep{WardKeller1993,IsaacsonNewby2013}. We
have not found the explicit formula in the literature, have not established its absence, and claim no priority.

\item \emph{What is proved about the mode (theorems, several computer-assisted;
Section~\ref{s6b_rigorous:ssec-status}).} Section~\ref{s6b_rigorous:sec} proves, for corner-to-corner transport (end to
end for $\dd=1$) and, where stated, for the other point-target placements: that the law is unimodal for the three
point-target placements studied in continuous time in every dimension and, in discrete time, for $q\le\frac45$ from
corner to corner in every dimension (for the centre placements when $\dd\ge2$), a threshold attained on the chain with
$N=4$, while unimodality fails near $q=1$ on the chain, at $q=1$ for some sizes in two dimensions, and for other
placements; that the residues of the poles do not alternate in sign for $\dd\ge2$, $N\ge3$; that in
one dimension $\lim\tmode/\MF$ exists for every
$q$ and is enclosed to within $10^{-70}$, with windows of width about one step for every $N\ge3$ when $q\le\frac45$;
and that in two and three dimensions the closed form has a scaled error $X\muone\,|\tmodec-q\tcf|$ bounded for every
$N\ge12$ ($\dd=2$) and $N\ge10$ ($\dd=3$), where $\tcf$ is the closed form \eqref{s6_mechanism:eq-L1-simple}, and
the leading laws $\tmode\sim(4/\pi^{2}q)N^{2}\ln\ln N$ and $(6/\pi^{2}q)N^{2}\ln N$ hold with explicit remainders ---
in continuous time and, in discrete time, for $q\le0.99$ once $N$ exceeds an explicit size and for every $q<1$ once $N$ exceeds an explicit size of order $\pi q/(1-q)$ as $q\to1$ (the case $q=1$ is
open) --- and that the closed form is within $1\%$ ($\dd=2$) and $0.3\%$ ($\dd=3$) of the certified exact mode on
finite ranges with $N\ge10$ for $q\in\{\frac45,\frac12\}$, in three dimensions for every $N\ge10$.
Computer-assisted proofs reduce a statement, by lemmas proved by hand, to finitely many inequalities checked in exact
or ball arithmetic; their trusted bases and certificates are named. Some ingredients are classical: the interlacing of
the two spectra \citep{HartichGodec2018,HartichGodec2019}, the mixture and exponential laws from the uniform and the
quasi-stationary start \citep{AldousFill}, and unimodality on the chain in continuous time \citep{Keilson1971}. We have not
found the other statements in the literature, but our search was targeted, not exhaustive; the tools (interlacing,
log-concavity under convolution, cone certificates) are classical.
% src: data/msc_rigorous_1d/03_constants.json; data/msc_rigorous_certified/closedform_summary_c128.json; data/article/s6b_rigorous_checks.json (corollary_closed_form_all_N)

\item \emph{Shape of the distribution (exact computation).} For $\dd\ge2$ the law is close to an exponential
preceded by a short dead time, whose end the mode marks; the mode is then not a typical time, and in three
dimensions the maximum is very flat (Section~\ref{s6_mechanism:ssec-shape}). This two-scale structure is known
\citep{Benichou2010,Meyer2011,GodecMetzler2016,AldousBrown1992}; we add the exact lattice computation.

\item \emph{Inert defects (exact computation for each configuration; dilute law formal).} Each configuration of
blocked sites on the square lattice is solved exactly, so the only uncertainty is the sampling of placements.
At $N=35$ and blocked fractions $p\le0.2$, uniformly random defects multiply the mean and the mode by nearly the same
factor $\Dzero/\Deff(p)$ (within $2.8\%$ and $1.82\%$, ensemble-mean estimates with $95\%$ intervals up to $\pm1.8\%$
and $\pm0.72\%$), so that $\tmode/\MF$ changes by a few per cent at low density but rises by $4.7\%$ to $8.9\%$ at
$p=0.2$, and where the defects sit decides the direction of any change of $\tmode/\MF$ (Section~\ref{s7_defects:sec}); a
single barrier on a chain already shows that placement decides which of the two times responds
\citep{SarvaharmanGiuggioli2023}.
% src: data/article/s7_defects_density_pooled.json (statements.uniform_max_abs_mfpt_fac_over_theta_pct_p_le_0.20_pooled, uniform_max_abs_mode_fac_over_theta_pct_p_le_0.20_pooled, uniform_ratio_excess_pct_by_p, uniform_rom_excess_pct_by_p); data/article/s7_defects_density_pooled.csv (mfpt_fac_ci_P, mode_fac_ci_P, theta_L192)
The dilute law
$\Deff/\Dzero=1-(\pi-1)p+\Order(p^{2})$ is in the literature \citep{WatsonLeath1974,Nieuwenhuizen1986,Ernst1987};
we prove the single-defect response behind it (Proposition~\ref{s7_defects_density:prop-dilute}). That the dilation
becomes exact as $N\to\infty$ is stated only as Conjecture~\ref{s7_defects_density:conj-homog}.
\end{enumerate}
\endgroup

\paragraph{Outline.}
Section~\ref{s2_model:sec} defines the model and its exact structure; Section~\ref{s3_methods:sec} describes the
computations and the certificate of the global mode; Section~\ref{s4_mfpt:sec} gives the exact means.
Sections~\ref{s5_modes:sec} and~\ref{s6_mechanism:sec} report the exact modes and explain them through the poles of
the first-passage transform, and Section~\ref{s6b_rigorous:sec} states what is proved about them.
Section~\ref{s7_defects:sec} treats inert defects, and Section~\ref{s9_discussion:sec} separates what is established
from what is not. The Supplementary Material (Sections~S1--S10) contains the proofs that are omitted here,
the extended tables, the validation of the numerical methods and the description of the certificates; the companion
notes R0--R5 in the repository \citep{ZhouyiRepo} contain the complete proofs of the computer-assisted theorems.
